%% ma: L10-B14-0086
%% lop: 10
%% chuong: 5
%% bai: 14
%% dang: 10.14.12
%% loai: DS
%% muc: [NB, NB, TH, VD]
%% dap_an: "ĐSĐĐ"
%% nguon: "chuyen-de-cac-so-dac-trung-cua-mau-so-lieu-khong-ghep-nhom-toan-10-chuong-trinh-moi.pdf (GV. Phan Nhật Linh) (P05), Bài 03 Dạng 2 Phần II Câu 3"
%% kien_thuc: "Lớp 10 Bài 13 (số trung bình, trung vị, tứ phân vị, mốt), Bài 14 (khoảng biến thiên, khoảng tứ phân vị, phương sai, độ lệch chuẩn, giá trị bất thường)"
%% trang_thai: chinh-thuc
%% ghi_chu: "Tự duyệt theo yêu cầu giáo viên 22/10/2026."
\begin{ex}
Mẫu số liệu sau cho biết số ghế trống tại một rạp chiếu phim trong $9$ ngày:
\begin{center}
\begin{tabular}{|c|c|c|c|c|c|c|c|c|}\hline
$7$ & $8$ & $22$ & $20$ & $18$ & $15$ & $19$ & $13$ & $11$\\\hline
\end{tabular}
\end{center}
Xét tính đúng sai của các khẳng định sau:
\choiceTF{\True Khoảng biến thiên của mẫu số liệu là $R=15$}{Trung vị của mẫu số liệu là $Q_2=18$}{\True Khoảng tứ phân vị của mẫu số liệu là $\Delta_Q=10$}{\True Phương sai của mẫu số liệu là $s^2\approx25{,}73$}
\loigiai{Sắp xếp: $7$; $8$; $11$; $13$; $15$; $18$; $19$; $20$; $22$.

a) Đúng: $R=22-7=15$.

b) Sai: $Q_2=15$.

c) Đúng: $Q_1=\dfrac{8+11}2=9{,}5$; $Q_3=\dfrac{19+20}2=19{,}5$; $\Delta_Q=10$.

d) Đúng: $\overline x=\dfrac{133}9\approx14{,}78$; $s^2\approx25{,}73$.}
\end{ex}
